\begin{center}
\emph{\S\ 9. Formenreihen. Reduktionss\"atze}\srcfn{*)}{Vgl. die entsprechenden Betrachtungen im tern\"aren Gebiete. Diss. \S\S\ 1--3.}.
\end{center}
Die in \S\ 6 (Schlu\ss) eingef\"uhrte \emph{Formenreihe} definieren wir jetzt etwas sch\"arfer als ,,die Gesamtheit der in den Koeffizienten linearen invarianten Bildungen einer gegebenen Ausgangsform, d. h. die Gesamtheit der durch Faltung in sich aus der Ausgangsform entstandenen Formen, \emph{in der Anordnung nach h\"oheren Formen}``.

Zur Erkl\"arung der h\"oheren Formen sei die Ausgangsform nach Satz VII als Produkt zweier Normalformen $S\cdot T$ gegeben; als \emph{h\"ohere Form} bezeichnen wir dann Formen mit h\"oherer Faltung in sich, unter den Formen mit gleicher Faltung in sich speziell normierte. Genauer ausgedr\"uckt hei\ss t das: Es sei die Form $A$ durch $\alpha_\rho$, die Form $B$ durch $\beta_\rho$ Grundfaltungen vom Typus $\rho$ entstanden. $B$ ist h\"oher als $A$:
\begin{enumerate}
\renewcommand{\labelenumi}{\arabic{enumi})}
\item wenn in dem System von Ungleichungen: $\beta_\rho\geqq\alpha_\rho$, $\rho=1,2,\ldots,n-1$, mindestens f\"ur einen Wert $\rho$ das Ungleichheitszeichen gilt\srcfn{**)}{Gilt in dem System von Ungleichungen zum Teil das Zeichen $>$, zum Teil das Zeichen $<$, so sind die Formen nicht miteinander vergleichbar, da einer Form, die aus einer anderen durch Faltung hervorgegangen, stets ein positives Wertesystem von Grundfaltungen zukommt, also in unserem Fall positive oder verschwindende Werte $\beta_\rho-\alpha_\rho$.};
\item wenn f\"ur alle Werte $\rho$ das Gleichheitszeichen gilt, aber weniger Symbolreihen zu einem einzigen Matrizenprodukt zusammengefa\ss t sind. Diese Normierung erweist sich als notwendig wegen der Reduktionen in \S\ 8. Danach wird z. B.: $\pair{S^{n-1}}{T_{n-1}}$ h\"oher als $\pair{S^{n-1}}{[S^1T^\rho]_{n-\rho-1}p_\rho}$.
\item Wenn f\"ur alle Werte $\rho$ das Gleichheitszeichen gilt und die Anzahl der zu einem Matrizenprodukt zusammengefa\ss ten Symbolreihen die gleiche ist, wird $B$ h\"oher als $A$ durch eine an sich willk\"urliche, aber ein f\"ur allemal festgelegte Normierung. Diese Normierung ist wegen der endlichen Anzahl der zu einem Wertesystem $\alpha$ geh\"orenden Formen stets m\"oglich, l\"a\ss t sich z. B. erreichen durch Auszeichnung der einen Form $S$ (gr\"o\ss ere Anzahl von Symbolreihen $S$, h\"ohere Summe der oberen Gewichtsindizes der Reihen $S$ usf.).
\end{enumerate}
Wir k\"onnen uns zweckm\"a\ss ig die Formenreihe als $(n-1)$-dimensionale Mannigfaltigkeit von Gitterpunkten angeordnet denken, wobei die $n-1$ Richtungen den $n-1$ Grundfaltungen entsprechen. Einer jeden Form ist dann ein und nur ein Wertesystem $\alpha$, d. h. ein positiv ganzzahliger Gitterpunkt zugeordnet, w\"ahrend die verschiedenen Formen, die zu einem Wertesystem $\alpha$ geh\"oren, durch die obige Normierung in ihrer Anordnung eindeutig bestimmt sind. Eine beliebige Form der Formenreihe gibt die Ausgangsform einer neuen Formenreihe, die als Teil in der Gesamtformenreihe enthalten ist; und zwar als Teil enthalten in denjenigen Formen, die man von der neuen Ausgangsform aus, in den $n-1$ positiven Richtungen fortschreitend, erh\"alt.

Verm\"oge Satz VIII und der allgemeinen Theorie der Reihenentwicklungen gelten nun f\"ur Formenreihen genau die Reduktionss\"atze des bin\"aren und tern\"aren Gebiets (vgl. Diss. \S\ 3), von denen wir den Hauptsatz noch kurz ableiten wollen. Wir schicken die Definition voraus:

,,In einem gegebenen System von Formen sei eine gewisse endliche Anzahl von Formen zu einem Modul zusammengefa\ss t; als \emph{reduzibel} bezeichnen wir eine Form, die sich ausdr\"ucken l\"a\ss t durch Formen, die Invarianten zum Faktor haben, oder durch ,,h\"ohere Formen``; d. h. durch h\"ohere Formen der Gesamtformenreihe, der die Form angeh\"ort, oder durch Formen, die die Symbole des Moduls in h\"oherer Ordnung enthalten. Eine reduzible Formenreihe nennen wir \emph{Reduzent}.`` Es gilt dann

\emph{Satz IX: Ist die Ausgangsform einer Formenreihe reduzibel dadurch, da\ss\ sie durch Faltung mit einem Reduzenten hervorgegangen ist, so ist die gesamte, aus dieser Ausgangsform entstehende Formenreihe reduzibel.}

Zum Beweise bedenken wir, da\ss\ eine durch Faltung einer Form $f$ mit einer zweiten $g$ entstandene Form $h$ sich wegen (45.) auffassen l\"a\ss t als eine mehrere kogrediente Variabelnreihen $p_\rho,p'_\rho$ enthaltende Form $f$. Solche Formen lassen sich aber nach der allgemeinen Theorie der Reihenentwicklung darstellen als Polaren der Elementarkovarianten von $f$, dadurch da\ss\ man die Reihenentwicklung sukzessive auf je zwei kogrediente Variabelnreihen $p_\rho,p'_\rho$ anwendet. Dabei sind nach (38.) die auftretenden Elementarkovarianten die durch kogrediente Faltung erzeugten Formen. Da aber die Reihenfolge in der Anwendung der Reihenentwicklung noch willk\"urlich ist, k\"onnen wir insbesondere eine solche Reihenfolge nehmen, die der Erzeugung der allgemeinen Faltungen aus Grundfaltungen entspricht. Das entstehende System von Elementarkovarianten wird dann identisch mit der Formenreihe $f$; und zwar ordnen sich die durch kogrediente Faltung von h\"oherem Defekte entstandenen Formen in die durch Grundfaltungen entstandenen ein. Zu Polaren der Formenreihe $f$ gelangt man aber auch bei beliebiger Reihenfolge der Reihenentwicklung, der ein zweites System von Elementarkovarianten entsprechen m\"oge. Da n\"amlich die Formenreihe die Gesamtheit der invarianten Bildungen ersch\"opft, l\"a\ss t sich jede Form des zweiten Systems linear ausdr\"ucken durch Polaren der Formenreihe, und zwar Polaren derjenigen Formen, die gleiche oder h\"ohere Gewichtsabnahme zeigen. Dies letztere folgt daraus, da\ss\ sich (vgl. \S\ 7) die Polaren auffassen lassen als Simultaninvarianten einer Formenreihe $H$ (52.), mit der Form $h$ entsprechenden Gradzahlen, und der Formenreihe $f$. Jeder Faltung von $H$ in sich entspricht eine Gewichtsabnahme, die sich nach Vornahme von Substitution (11.) auf die Symbolreihen und damit auf die Formen von $f$ \"ubertr\"agt\srcfn{*)}{Diese Betrachtung liegt auch der zweiten Fassung des \"Aquivalenzbegriffes (\S\ 8 Schlu\ss) zugrunde. Weitere Ausf\"uhrungen \"uber die verschiedenen Systeme von Elementarkovarianten finden sich im tern\"aren Gebiet bei Study, a. a. O., II. \S\ 7. Satz 7) und 8). Verm\"oge unseres Satz VII gelten dieselben auch f\"ur $n$ Variable.}.

Eine beliebige Form der Formenreihe $h$ ist entstanden durch Faltung von $h$ in sich; d. h. durch h\"ohere Faltung von $g$ mit $f$ einerseits, durch Faltung von $f$ oder $g$ in sich andererseits. In beiden F\"allen f\"uhrt die Reihenentwicklung, ebenso wie oben bei der Ausgangsform $h$ gezeigt, auf Polaren der Formenreihe $f$. Ist insbesondere $f$ Reduzent, so ergibt sich eine Entwicklung nach Polaren des Reduzenten; die Formenreihe $h$ wird also reduzibel.

Anwendungen des Satzes auf die Reduktion vollst\"andiger Formensysteme ergeben sich analog wie im bin\"aren und tern\"aren Gebiet.
